% © 2026 Ing. David Jaroš — CC BY-NC-ND 4.0
%
% This work is licensed under a Creative Commons Attribution-NonCommercial-NoDerivatives
% 4.0 International License (CC BY-NC-ND 4.0).
%
% License History: Earlier drafts (up to v0.3) were released under CC BY 4.0.
% From v0.4 onward, all material is released under CC BY-NC-ND 4.0 to protect
% the integrity of the theoretical work during ongoing academic development.

% UBT-AI-PROVENANCE-BEGIN
% schema: ubt-ai-provenance/v1
% tier: C_working
% ai_assistance: disclosed
% human_review: risk-based
% editorial_responsibility: Ing. David Jaroš
% policy: ../../AI_PROVENANCE.md
% notice: Working material; exhaustive human review is not claimed.
% UBT-AI-PROVENANCE-END

\documentclass[11pt,a4paper]{article}
\usepackage{amsmath,amssymb,amsthm}
\usepackage{booktabs}
\usepackage{geometry}
\usepackage{hyperref}
\geometry{margin=2.5cm}

\newtheorem{proposition}{Proposition}[section]
\newtheorem{remark}[proposition]{Remark}
\newtheorem{lemma}[proposition]{Lemma}

\title{Parent Coupling $e_5$ Audit in UBT\\
\large Target 3 — fix\_or\_reject\_U1\_coupling\_normalization}
\author{Ing.\ David Jaroš}
\date{2026-05-10}

\begin{document}
\maketitle

\begin{abstract}
We audit whether the higher-dimensional electromagnetic coupling $e_5$,
appearing in the 5D covariant derivative $D_M=\partial_M+ie_5 A_M$, is fixed
by the canonical UBT action $S[\Theta]$.  We check field-rescaling freedom,
canonical normalization, curvature and trace normalization constraints,
topological quantization, and Chern-class winding conditions.
Verdict: \textbf{NO-GO}.
\end{abstract}

% ----------------------------------------------------------------
\section{Where $e_5$ enters the UBT action}
% ----------------------------------------------------------------

The five-dimensional sector of $S[\Theta]$ relevant to the electromagnetic
coupling is
\begin{equation}
  S[\Theta]\supset
  \int d^5x\;\mathrm{Tr}\!\left[(D_M\Theta)^\dagger(D^M\Theta)\right]
  -\frac{1}{4e_5^2}\int d^5x\;F_{MN}F^{MN},
\end{equation}
where $M=(\mu,\psi)$ runs over all five dimensions,
\begin{equation}
  D_M\Theta = \partial_M\Theta + ie_5 A_M\Theta,
\end{equation}
and $F_{MN}=\partial_M A_N - \partial_N A_M$.

The coupling $e_5$ has mass dimension $[e_5]=[\text{mass}]^{-1/2}$ in 5D
(mass dimension analysis: $[A_M]=[\text{mass}]^{3/2}$,
$[D_M\Theta]=[D_M][\Theta]$, $[e_5][\Theta][A_M]=[\partial_M][\Theta]$).

% ----------------------------------------------------------------
\section{Field-rescaling freedom}
% ----------------------------------------------------------------

\begin{lemma}[Absorption of $e_5$ by field rescaling]
The coupling $e_5$ can always be absorbed into a redefinition of $A_M$:
\begin{equation}
  \mathcal{A}_M := e_5 A_M,
  \qquad
  D_M = \partial_M + i\mathcal{A}_M.
\end{equation}
In terms of $\mathcal{A}_M$, the matter kinetic term is $e_5$-free:
\begin{equation}
  \mathrm{Tr}[(D_M\Theta)^\dagger(D^M\Theta)]
  =\mathrm{Tr}\!\left[(\partial_M\Theta+i\mathcal{A}_M\Theta)^\dagger
    (\partial^M\Theta+i\mathcal{A}^M\Theta)\right].
\end{equation}
The gauge kinetic term becomes
\begin{equation}
  -\frac{1}{4e_5^2}F_{MN}F^{MN}
  =-\frac{1}{4}\mathcal{F}_{MN}\mathcal{F}^{MN},
\end{equation}
with $\mathcal{F}_{MN}=\partial_M\mathcal{A}_N-\partial_N\mathcal{A}_M$.
\end{lemma}

In the rescaled basis, the entire action is independent of $e_5$.  Therefore,
$e_5$ is a pure normalization convention in the absence of external constraints.

\textbf{Status: NO-GO (field-rescaling removes $e_5$).}

% ----------------------------------------------------------------
\section{Canonical normalization of $\Theta$}
% ----------------------------------------------------------------

Canonical normalization of the kinetic term requires the matter part of the
action to take the form
\begin{equation}
  S_\Theta = \int d^5x\;\mathrm{Tr}\!\left[(\partial_M\Theta)^\dagger(\partial^M\Theta)\right]
  + \text{interaction terms}.
\end{equation}
This fixes the normalization of $\Theta$ relative to $\partial_M$, but places
no constraint on the relative coupling between $\Theta$ and $A_M$ beyond what
gauge invariance requires.  The coupling constant $e_5$ remains a free
multiplicative factor in the interaction vertex.

\textbf{Status: NO-GO (canonical normalization of $\Theta$ does not fix $e_5$).}

% ----------------------------------------------------------------
\section{Curvature and trace normalization}
% ----------------------------------------------------------------

The Dirac quantization condition for a U(1) gauge theory on a manifold with
boundary requires the first Chern class to satisfy
\begin{equation}
  c_1 = \frac{1}{2\pi}\int_\Sigma F_{MN}dS^{MN} \in \mathbb{Z}.
\end{equation}
This is a topological condition on the \emph{integrated} flux and constrains
the product $e_5\cdot\Phi$, not $e_5$ alone.  If the magnetic flux $\Phi$ is
unconstrained, $e_5$ remains free.

\begin{remark}[Trace normalization in the gauge sector]
The gauge kinetic term in the non-abelian sector uses
$\mathrm{Tr}(T^aT^b)=\tfrac{1}{2}\delta^{ab}$.  For the U(1) factor, the
normalization is $\mathrm{Tr}(T_{\mathrm{EM}}^2)=1/2$ (derived in Target 1).
This fixes the \emph{relative} normalization between the U(1) and SU(2) kinetic
terms, but not the overall coefficient $1/e_5^2$ in front of the gauge kinetic
action.
\end{remark}

\textbf{Status: NO-GO (Chern class fixes integer flux, not $e_5$; trace
normalization fixes relative coefficient, not absolute coupling).}

% ----------------------------------------------------------------
\section{Topological quantization and Chern-class winding}
% ----------------------------------------------------------------

On the compactification circle $S^1_\psi$, the holonomy
\begin{equation}
  \mathcal{W}=\exp\!\left(ie_5\oint_{S^1_\psi}A_\psi\,d\psi\right)
\end{equation}
must equal $\pm 1$ (or more generally an element of the stabilizer group of the
vacuum) for single-valuedness of $\Theta$.  This gives
\begin{equation}
  e_5\oint A_\psi\,d\psi = 2\pi n,\qquad n\in\mathbb{Z}.
\end{equation}
The integer $n$ is the winding number.  For the minimal case $n=1$,
\begin{equation}
  e_5 = \frac{2\pi}{\oint A_\psi\,d\psi}.
\end{equation}
This fixes $e_5$ only if the flux $\oint A_\psi\,d\psi$ is independently
fixed.  In the abelian gauge theory on $S^1_\psi$, the flux is a modulus
(Wilson line modulus) that is not constrained by the action.

\textbf{Status: NO-GO (winding quantization relates $e_5$ to an unconstrained
flux modulus).}

% ----------------------------------------------------------------
\section{Flux quantization}
% ----------------------------------------------------------------

In a magnetic compactification, flux quantization on a two-cycle $\Sigma_2$
gives
\begin{equation}
  \frac{e_5}{2\pi}\int_{\Sigma_2} F = n \in \mathbb{Z}.
\end{equation}
On the torus $T^2=S^1_t\times S^1_\psi$, this fixes
$e_5 \cdot B \cdot (2\pi R_t)(2\pi R_\psi) \in 2\pi\mathbb{Z}$.  Again,
the product $e_5 B R_t R_\psi$ is quantized, not $e_5$ individually.

\textbf{Status: NO-GO (flux quantization relates $e_5$ to unconstrained flux
density and radii).}

% ----------------------------------------------------------------
\section{Summary and verdict}
% ----------------------------------------------------------------

\begin{center}
\begin{tabular}{@{}lll@{}}
\toprule
Check & Outcome & Reason \\
\midrule
Field rescaling $\mathcal{A}_M=e_5 A_M$ & NO-GO & $e_5$ absorbed, action $e_5$-free \\
Canonical $\Theta$ normalization & NO-GO & Does not constrain coupling \\
Chern class quantization & NO-GO & Quantizes $e_5\cdot\Phi$, not $e_5$ \\
Trace normalization in gauge sector & NO-GO & Fixes relative, not absolute \\
Holonomy / winding on $S^1_\psi$ & NO-GO & Relates $e_5$ to Wilson line modulus \\
Flux quantization on $T^2$ & NO-GO & Fixes $e_5 B R^2$ product only \\
\bottomrule
\end{tabular}
\end{center}

\paragraph{Overall verdict.}
\textbf{NO-GO.}
The coupling $e_5$ is a free parameter in the current UBT formulation.  No
mechanism in $S[\Theta]$, topological quantization conditions, or trace
normalization requirements uniquely fixes the value of $e_5$.  Field rescaling
shows that $e_5$ can always be removed from the action by a redefinition of
$A_M$, confirming it is a pure normalization convention.

\begin{remark}[Possible route to fixing $e_5$]
A non-perturbative constraint relating the gauge sector normalization to the
gravitational sector (e.g., through a Planck-scale relation
$e_5^2\sim G_5\cdot\ell_s^{-3}$) could in principle fix $e_5$.  This would
require a unification condition linking the gravitational and electromagnetic
sectors of $S[\Theta]$, which is not present in the current canonical
formulation.
\end{remark}

\end{document}
